\documentclass{article}
\usepackage{amsmath,amssymb,graphicx}
\usepackage{geometry}
\geometry{margin=2.5cm}

\begin{document}

\title{Evaluation of increment PDFs using Laplace's method}
\author{}
\date{}
\maketitle

\begin{abstract}
We study probability density functions of Eulerian velocity increments in turbulent flows using the multifractal propagator formalism with the Kolmogorov--Oboukhov lognormal model. The increment PDFs are expressed as integrals over the Hölder exponent that lack closed-form solutions. We apply Laplace's method to derive asymptotic approximations in the inertial range, examining two saddle-point approaches: one assuming a fixed dominant exponent and another using the Lambert W function. Numerical comparisons assess the validity of each approximation and provide a basis for entropy-based intermittency analyses.
\end{abstract}

\section{Introduction}

Turbulent flows exhibit intermittency: velocity fluctuations become increasingly non-Gaussian at smaller scales. Understanding the statistics of velocity increments $\delta_\ell u = u(x+\ell) - u(x)$ across scales $\ell$ is fundamental to turbulence theory.

The multifractal approach characterizes intermittency through a distribution of local Hölder exponents $h$ with a scale-dependent spectrum $D_E(h)$. Chevillard et al.\ \cite{chevillard2012} developed a propagator formalism coupling large-scale Gaussian fluctuations with small-scale intermittency, expressing increment PDFs as integrals over the Hölder distribution.

We adopt the Kolmogorov--Oboukhov lognormal model \cite{kolmogorov1962, obukhov1962}, where $D_E(h) = 1 - (h-c_1)^2/(2c_2)$. This model provides analytical tractability while capturing key intermittency features. The resulting propagator integrals cannot be evaluated exactly, motivating our use of Laplace's method for asymptotic analysis in the limit of small scale separation.

The paper is organized as follows. Section 2 derives the Hölder distribution and propagator integral. Section 3 presents two Laplace approximations: fixed saddle point (Case A) and exact solution via Lambert W (Case B). Section 4 introduces a normalized PDF based on second-order structure functions. Section 5 discusses entropy calculations. Numerical results compare both approximations with direct integration, clarifying their validity ranges.

\section{Starting point: the distribution of $h$}

We start from (Chevillard et al., 2012, Eq.~11), which defines the probability density function of the local Hölder exponent \(h\):
\begin{equation}
\mathcal{P}_h^{(\ell)}(h)
=
\frac{
\left(\tfrac{\ell}{L}\right)^{1-\mathcal{D}^E(h)}
}{
\displaystyle \int_{h_{\min}}^{h_{\max}}
\left(\tfrac{\ell}{L}\right)^{1-\mathcal{D}^E(h)}\,\mathrm{d}h
}.
\tag{1}
\end{equation}

using the Kolmogorov–Oboukhov log-normal spectrum
\[
\mathcal{D}^E(h)=1-\frac{(h-c_1)^2}{2c_2},
\]
with \(h_{\min}=-\infty\) and \(h_{\max}=+\infty\),
we obtain
\begin{equation}
\mathcal{P}_h^{(\ell)}(h)
=
\frac{
\exp\!\left[-\tfrac{(h - c_1)^2}{2c_2}\ln\!\left(\tfrac{L}{\ell}\right)\right]
}{
\displaystyle \int_{-\infty}^{+\infty}
\exp\!\left[-\tfrac{(h - c_1)^2}{2c_2}\ln\!\left(\tfrac{L}{\ell}\right)\right]
\mathrm{d}h
}.
\tag{2}
\end{equation}

To simplify notation, we define

\begin{equation}
\alpha = -\ln\!\left(\frac{\ell}{L}\right) > 0.
\tag{3}
\end{equation}

Then the expression becomes
\begin{equation}
\mathcal{P}_h^{(\ell)}(h)
=
\frac{
\exp\!\left[-\frac{\alpha}{2c_2}(h - c_1)^2\right]
}{
\displaystyle \int_{-\infty}^{+\infty}
\exp\!\left[-\frac{\alpha}{2c_2}(h - c_1)^2\right]
\mathrm{d}h
}.
\tag{4}
\end{equation}

The denominator is a Gaussian integral:
\begin{equation}
\int_{-\infty}^{+\infty}
e^{-\frac{\alpha}{2c_2}(h - c_1)^2}\mathrm{d}h
= \sqrt{\frac{2\pi c_2}{\alpha}}.
\tag{5}
\end{equation}

Hence,
\begin{equation}
\boxed{
\mathcal{P}_h^{(\ell)}(h)
=
\sqrt{\frac{\alpha}{2\pi c_2}}\,
\exp\!\left[-\frac{\alpha}{2c_2}(h - c_1)^2\right].
}
\tag{6}
\end{equation}

\section{Computation of the PDF using Laplace's method}

\subsection{General form of the PDF}

We now consider the probability density function of velocity increments. According to (Chevillard et al., 2012, Eq.~17) :
\begin{equation}
\mathcal{P}_{\delta_\ell u}(x)
= \frac{1}{2\pi\sigma}\,
\sqrt{\frac{\alpha}{c_2}}
\int_{-\infty}^{+\infty}
\exp[-f(h,\alpha,x)]\,\mathrm{d}h.
\tag{7}
\end{equation}

with
\begin{equation}
f(h,\alpha,x)
= \alpha\!\left(\frac{(h-c_1)^2}{2c_2}-h\right)
+ \frac{1}{2}\left(\frac{x}{\sigma}\right)^{2}
e^{2\alpha h}
\tag{8}
\end{equation}



$\\[0.8em]$
We would like to get an explicit form of our Eq.~7. To do so, we'll use Laplace's approximation (see Appendix).
Knowing that
\begin{equation}
\forall h \in \mathbf{R},
\quad
\frac{\partial f}{\partial h}(h,\alpha,x)
= \alpha\!\left(\frac{h-c_1}{c_2}-1\right)
+ \alpha\!\left(\frac{x}{\sigma}\right)^{2}
e^{2\alpha h}
\tag{9}
\end{equation}

we have :
\begin{equation}
\forall h \in \mathbf{R},
\quad
\frac{\partial^2 f}{\partial h^2}(h,\alpha,x)
=
\frac{\alpha}{c_2}
+ 2\alpha^2
\left(\frac{x}{\sigma}\right)^{2}
e^{2\alpha h}
\quad > 0
\tag{10}
\end{equation}

We just need to verify that $h^\star$ is such that
$\frac{\partial f}{\partial h}(h^\star,\alpha,x) = 0$
to be able to use Laplace approximation.

$\\[0.3em]$
In this case, by applying the approximation to Eq.~7, we get :

\begin{equation}
\boxed{
\mathcal{P}_{\delta_\ell u}(x)
\approx
\sqrt{\frac{\alpha}{
2\pi c_2 \sigma^2 \frac{\partial^2 f}{\partial h^2}(h,\alpha,x))
}}\,
\exp[-f(h^\star,\alpha,x)].
}
\tag{11}
\end{equation}


$\\[0.1em]$
i.e
\begin{equation}
\boxed{
\begin{aligned}
\mathcal{P}_{\delta_\ell u}(x)
\approx
\frac{1}
{
\sqrt{2\pi}
\sqrt{
2\alpha c_2 x^2 e^{2\alpha h^{\star}}
+ \sigma^2}}
\times
\exp\!\left[
-\alpha\!\left(
\frac{(h^\star-c_1)^2}{2c_2}-h^\star
\right)
-\frac{1}{2}
\left(\frac{x}{\sigma}\right)^{2}
e^{2\alpha h^\star}
\right]
\end{aligned}
}
\tag{12}
\end{equation}


\subsection{Case A: assuming \(h^\star = c_1\)}

As a first approximation, we can consider that :

\begin{equation}
\frac{\partial f}{\partial h}(h^\star,\alpha,x)\Big|_{h^\star=c_1}
= 0
\tag{13}
\end{equation}

$\\[0.1em]$
We can use the Laplace approximation given by Eq.~12, which becomes :

\begin{equation}
\boxed{
\begin{aligned}
\mathcal{P}_{\delta_\ell u}(x)
\approx
\frac{1}{\sqrt{2\pi}}
\frac{1}
{
\sqrt{
2\alpha c_2 x^2 e^{2\alpha c_{1}}
+ \sigma^{2}}}
\times
\exp\!\left[
\alpha c_1
-\frac{1}{2}\left(\frac{x}{\sigma}\right)^{2}
e^{2\alpha c_1}
\right]
\end{aligned}
}
\tag{14}
\end{equation}

\subsection{Case B: solving \(\frac{\partial f}{\partial h}(h,\alpha,x)=0\) via the Lambert \(W\) function}

From \((8)\),
\begin{equation}
\frac{\partial f}{\partial h}(h,\alpha,x)
= \alpha\!\left(\frac{h-c_1}{c_2}-1\right)
+ \alpha\!\left(\frac{x}{\sigma}\right)^{2}
e^{2\alpha h}
\tag{15}
\end{equation}
Set to zero:
\begin{equation}
h - (c_1 + c_2)
= -c_2
\left(\frac{x}{\sigma}\right)^{2}
e^{2\alpha h}
\tag{16}
\end{equation}
Multiply by \(-2\alpha\) and let \(u=-2\alpha(h-c_1-c_2)\):
\begin{equation}
u
= 2\alpha c_2
\left(\frac{x}{\sigma}\right)^{2}
e^{2\alpha(c_1 + c_2)-u}
\tag{17}
\end{equation}
Consequently
\begin{equation}
u e^{u}
= 2\alpha c_2
\left(\frac{x}{\sigma}\right)^{2}
e^{2\alpha(c_1 + c_2)}
\tag{18}
\end{equation}
Hence
\begin{equation}
\boxed{
h^\star
= (c_1 + c_2)
- \frac{1}{2\alpha}
W\!\left(
2\alpha c_2
\left(\frac{x}{\sigma}\right)^{2}
e^{2\alpha(c_1 + c_2)}
\right)
}
\tag{19}
\end{equation}
where \(W(x)\) is the Lambert \(W\) function (the inverse of \(w\mapsto we^w\)).

$\\[0.1em]$
We can now plug the obtained $h^\star$ into Eq.~12.


\section{Normalized PDF}

According to (Chevillard et al., 2012, Eq.~12), the $p$-th order moment of the velocity increment
is expressed as
\begin{equation}
\langle |\delta_\ell u|^p \rangle \underset{\ell \to 0}{\approx} \sigma^2 \langle |\delta|^p \rangle \left(\frac{\ell}{L}\right)^{\zeta_p},
\quad\text{with}\quad
\zeta_p = \min_h \bigl[ p h + 1 - \mathcal{D}^E(h) \bigr] \quad\text{and}\quad \delta \sim \mathcal{N}(0,1).
\tag{20}
\end{equation}

Using the Kolmogorov–Oboukhov log-normal spectrum
\begin{equation}
\mathcal{D}^E(h) = 1 - \frac{(h - c_1)^2}{2c_2},
\tag{21}
\end{equation}

the minimization problem is
\begin{equation}
\frac{\mathrm{d}}{\mathrm{d}h}
\bigl[p h + 1 - \mathcal{D}^E(h)\bigr] = 0
\quad i.e. \quad
\frac{\mathrm{d}}{\mathrm{d}h}
\bigl[p h + \frac{(h-c_1)^2}{2c_2}\bigr] = 0
\tag{22}
\end{equation}

Hence, the minimizing exponent is
\begin{equation}
h_p = c_1 - p c_2,
\tag{23}
\end{equation}
and inserting this value into Eq.~20 gives
\begin{equation}
\zeta_p = p c_1 - \frac{p^2}{2} c_2
\tag{24}
\end{equation}

In particular, for the second-order :
\begin{equation}
\zeta_2 = 2(c_1 - c_2)
\tag{25}
\end{equation}

Knowing that $\langle |\delta|^2 \rangle = 1$, we get :
\begin{equation}
\boxed{
\langle (\delta_\ell u)^2 \rangle
\underset{\ell \to 0}{\approx}
\sigma^2 \left( \frac{\ell}{L} \right)^{2(c_1 - c_2)}.
}
\tag{26}
\end{equation}

We now define a normalized increment
\begin{equation}
\boxed{
\delta_\ell \nu
= \frac{\delta_\ell u}{\sqrt{\langle (\delta_\ell u)^2 \rangle}}
= \frac{\delta_\ell u}{\sigma \left( \frac{\ell}{L} \right)^{c_1 - c_2}}
}
\tag{27}
\end{equation}

$\\[0.3em]$
This normalization ensures that
\begin{equation}
\langle (\delta_\ell \nu)^2 \rangle = 1,
\tag{28}
\end{equation}
for all scales $\ell/L$, so that all normalized PDFs can be directly compared across scales.

$\\[0.1em]$
We thus have $\delta_\ell \nu\ = \delta_\ell u / \sigma \left( \frac{\ell}{L} \right)^{c_1 - c_2} = g(\delta_\ell u)$ where $g(x) = x / \sigma \left( \frac{\ell}{L} \right)^{c_1 - c_2}$ is monoton and differentiable.


$\\[0.1em]$
Knowing that : $\mathcal{P}_{\delta_\ell \nu}(x) = \left| \frac{1}{g'(g^{-1}(x))} \right| \mathcal{P}_{\delta_\ell u}(g^{-1}(x))$
, we have : $\mathcal{P}_{\delta_\ell \nu}(x) = \sigma \left( \frac{\ell}{L} \right)^{c_1 - c_2} \mathcal{P}_{\delta_\ell u}(\sigma \left( \frac{\ell}{L} \right)^{c_1 - c_2}x)$
i.e.

\begin{equation}
\boxed{
\mathcal{P}_{\delta_\ell \nu}(x) = \sigma e^{\alpha (c_2 - c_1)} \mathcal{P}_{\delta_\ell u}(\sigma e^{\alpha (c_2 - c_1)}x)
}
\tag{29}
\end{equation}

$\\[0.1em]$
We then obtain from Eq.~12 the normalized form of the PDF :

\begin{equation}
\boxed{
\begin{aligned}
\mathcal{P}_{\delta_\ell \nu}(x)
\approx
\frac{e^{\alpha(c_2 - c_1)}}{
\sqrt{2\pi}
\sqrt{
2 \alpha c_2 x^2 e^{2\alpha (h^\star + c_2 - c_1)} + 1
}
}
\times
\exp\!\left[
-\alpha\!\left(
\frac{(h^\star-c_1)^2}{2c_2}-h^\star
\right)
- \frac{1}{2}
x^{2}
e^{2\alpha (h^\star + c_2 - c_1)}
\right]
\end{aligned}
}
\tag{30}
\end{equation}

which is valid if $h^\star$ is such that $\frac{\partial f}{\partial h}(h^\star,\alpha,\sigma e^{\alpha (c_2-c_1)} x)=0$.


\subsection{Case A: assuming \(h^\star = c_1\)}

If we consider that
$\frac{\partial f}{\partial h}(h^\star,\alpha,\sigma e^{\alpha (c_2-c_1)} x)\Big|_{h^\star=c_1} = 0$,
then we get from Eq.~30 :

\begin{equation}
\boxed{
\begin{aligned}
\mathcal{P}_{\delta_\ell \nu}(x)
\approx
\frac{1}{
\sqrt{2\pi}
\sqrt{
2 \alpha c_2\,
x^2
e^{2\alpha c_2}
+ 1
}
}
\times
\exp\!\left[
\alpha c_2
-\frac{1}{2}
x^2
e^{2\alpha c_2}
\right]
\end{aligned}
}
\tag{31}
\end{equation}



\subsection{Case B: solving \(\frac{\partial f}{\partial h}(h,\alpha,\sigma e^{\alpha (c_2-c_1)} x)=0\) via the Lambert \(W\) function}

We want to find \(h^\star\) such that $\frac{\partial f}{\partial h}(h^\star,\alpha,\sigma e^{\alpha (c_2-c_1)}x) = 0$.
In Eq.~19, by considering that \(h^\star\) is a function of x, we just need to replace x by \( \sigma e^{\alpha (c_2-c_1)} x\). We get :

\begin{equation}
\boxed{
h^\star
= (c_1 + c_2)
- \frac{1}{2\alpha}
W\!\left(
2\alpha c_2 x^2 e^{4\alpha c_2}
\right)
}
\tag{32}
\end{equation}

We can now plug the obtained $h^\star$ into Eq.~30. Let $W=W\!\left( 2\alpha c_2 x^2 e^{4\alpha c_2} \right)$. We get :

\begin{equation}
\boxed{
\mathcal{P}_{\delta_\ell \nu}(x)
\approx
\frac{1}{
\sqrt{2\pi}\sqrt{2 \alpha c_2 x^2 e^{4 \alpha c_2-W} + 1}}
\times
\exp\left[
-\frac{\alpha}{2} \left( \frac{1}{c_2} \left( \frac{W}{2\alpha} \right)^2 - 3c_2 \right)
-\frac{1}{2} x^2 e^{4 \alpha c_2 - W}
\right]
}
\tag{33}
\end{equation}

i.e.

\begin{equation}
\boxed{
\mathcal{P}_{\delta_\ell \nu}(x)
\approx
\frac{e^\frac{3}{2}c_2 \alpha}{\sqrt{2\pi}}
\frac{e^{-\frac{2W + W^2}{8c_2 \alpha}}}{\sqrt{1+W}}
}
\tag{34}
\end{equation}




\section{Entropy of the normalized PDF}

We now compute the Shannon entropy of the normalized increment PDF
\(\mathcal{P}_{\delta_\ell \nu}\) obtained in Eq.~28. By definition,
\begin{equation}
S(\delta_\ell \nu)
=
- \int_{-\infty}^{+\infty}
\mathcal{P}_{\delta_\ell \nu}(x)\,
\ln\!\bigl[\mathcal{P}_{\delta_\ell \nu}(x)\bigr]\,
\mathrm{d}x
\tag{35}
\end{equation}




\section*{References}
\noindent
[1] L. Chevillard, B. Castaing, A. Arneodo, E. Lévêque, J.-F. Pinton, and S. G. Roux,
\emph{A phenomenological theory of Eulerian and Lagrangian velocity fluctuations in turbulent flows},
C.~R.~Physique, \textbf{13}, 899–928 (2012).

\appendix
\section{Appendix - Laplace's approximation theorem}
Let $I=[a,b]$ be an interval, $\varphi\in C^2(I,\mathbb{R})$
and $g\in C^0(I,\mathbb{C})$. Assume that:
\begin{enumerate}
    \item[(i)] For all $t>0$, $\displaystyle \int_a^b e^{t\varphi(x)} |g(x)|\,dx < +\infty$;
    \item[(ii)] $\varphi'$ vanishes at a unique point $x_0\in I$ and $\varphi''(x_0)<0$;
    \item[(iii)] $g(x_0)\neq 0$.
\end{enumerate}
Then
\[
\int_a^b e^{\,t\varphi(x)}\, g(x)\,dx
\sim
\sqrt{\frac{2\pi}{t\,|\varphi''(x_0)|}}\,
e^{\,t\varphi(x_0)}\, g(x_0)
\qquad\text{as } t\to +\infty .\]

\end{document}
